\documentclass[11pt]{article}
\usepackage[T1]{fontenc}
\usepackage{textcomp}
\usepackage[margin=0.75in]{geometry}
\usepackage{amsmath,amssymb,amsthm}
\usepackage{array,booktabs}
\usepackage{hyperref}
\setlength{\parindent}{0pt}
\newtheorem{theorem}{Theorem}
\theoremstyle{definition}
\newtheorem{definition}{Definition}
\title{Flow Proof Artifact: Finset}
\author{Generated by \texttt{flow doc proof}}
\date{Flow Proof Book}
\begin{document}
\maketitle
Finite set algebra on Finset carriers.\\[0.5em]
\textbf{Source.} graham-knuth-patashnik — *Concrete Mathematics*\\[0.5em]
\bigskip
\noindent{\large \textbf{Definition 1.} \textit{Union with the empty set changes nothing} \hfill Finset \cdot union \cdot empty is the right identity}\par
\smallskip\noindent\textit{Coordinate: Finset \cdot union \cdot empty is the right identity}\par
\smallskip\noindent\textit{Source: graham-knuth-patashnik}\par
\medskip
\noindent\textbf{Goal.} Union with the empty set changes nothing\par
\noindent$\displaystyle \forall s \in Finset\quad s ∪ empty = s$\par
\medskip
\noindent
\begin{tabular}{@{} >{\raggedright\arraybackslash}p{0.06\textwidth} >{\raggedright\arraybackslash}p{0.47\textwidth} >{\raggedleft\arraybackslash}p{0.41\textwidth} @{}}
\textbf{\#} & \textbf{Proof} & \textbf{Mathematics} \\
\midrule
\textbf{1.} & We stipulate empty is the right identity for union on Finset — this is a definition, not a derived fact. & \\[0.45em]
\textbf{2.} & This follows directly from the definition: s ∪ empty equals s. Hence proven. & $\displaystyle s ∪ empty = s$ \\[0.45em]
\bottomrule
\end{tabular}
\label{thm:Finset-union-empty_is_the_right_identity}
\par\medskip\hrule
\bigskip
\noindent{\large \textbf{Definition 2.} \textit{Union with the empty set on the left changes nothing} \hfill Finset \cdot union \cdot empty is the left identity}\par
\smallskip\noindent\textit{Coordinate: Finset \cdot union \cdot empty is the left identity}\par
\smallskip\noindent\textit{Source: graham-knuth-patashnik}\par
\medskip
\noindent\textbf{Goal.} Union with the empty set on the left changes nothing\par
\noindent$\displaystyle \forall s \in Finset\quad empty ∪ s = s$\par
\medskip
\noindent
\begin{tabular}{@{} >{\raggedright\arraybackslash}p{0.06\textwidth} >{\raggedright\arraybackslash}p{0.47\textwidth} >{\raggedleft\arraybackslash}p{0.41\textwidth} @{}}
\textbf{\#} & \textbf{Proof} & \textbf{Mathematics} \\
\midrule
\textbf{1.} & We stipulate empty is the left identity for union on Finset — this is a definition, not a derived fact. & \\[0.45em]
\textbf{2.} & This follows directly from the definition: empty ∪ s equals s. Hence proven. & $\displaystyle empty ∪ s = s$ \\[0.45em]
\bottomrule
\end{tabular}
\label{thm:Finset-union-empty_is_the_left_identity}
\par\medskip\hrule
\bigskip
\noindent{\large \textbf{Derived fact 3.} \textit{Union order does not matter} \hfill Finset \cdot union \cdot order does not matter}\par
\smallskip\noindent\textit{Coordinate: Finset \cdot union \cdot order does not matter}\par
\smallskip\noindent\textit{Source: graham-knuth-patashnik}\par
\medskip
\noindent\textbf{Goal.} Union order does not matter\par
\noindent$\displaystyle \forall a \in Finset \forall b \in Finset\quad a ∪ b = b ∪ a$\par
\medskip
\noindent
\begin{tabular}{@{} >{\raggedright\arraybackslash}p{0.06\textwidth} >{\raggedright\arraybackslash}p{0.47\textwidth} >{\raggedleft\arraybackslash}p{0.41\textwidth} @{}}
\textbf{\#} & \textbf{Proof} & \textbf{Mathematics} \\
\midrule
\textbf{1.} & We prove that order does not matter for union on Finset. & \\[0.45em]
\textbf{2.} & We can deduce that a ∪ b equals b ∪ a. Hence proven. & $\displaystyle a ∪ b = b ∪ a$ \\[0.45em]
\bottomrule
\end{tabular}
\label{thm:Finset-union-order_does_not_matter}
\par\medskip\hrule
\bigskip
\noindent{\large \textbf{Definition 4.} \textit{The empty finite set has cardinality zero} \hfill Finset \cdot cardinality \cdot the empty set has size zero}\par
\smallskip\noindent\textit{Coordinate: Finset \cdot cardinality \cdot the empty set has size zero}\par
\smallskip\noindent\textit{Source: graham-knuth-patashnik}\par
\medskip
\noindent\textbf{Goal.} The empty finite set has cardinality zero\par
\noindent$\displaystyle card(empty) = 0$\par
\medskip
\noindent
\begin{tabular}{@{} >{\raggedright\arraybackslash}p{0.06\textwidth} >{\raggedright\arraybackslash}p{0.47\textwidth} >{\raggedleft\arraybackslash}p{0.41\textwidth} @{}}
\textbf{\#} & \textbf{Proof} & \textbf{Mathematics} \\
\midrule
\textbf{1.} & We stipulate the empty set has size zero for cardinality on Finset — this is a definition, not a derived fact. & \\[0.45em]
\textbf{2.} & This follows directly from the definition: card(empty) equals 0. Hence proven. & $\displaystyle card(empty) = 0$ \\[0.45em]
\bottomrule
\end{tabular}
\label{thm:Finset-cardinality-the_empty_set_has_size_zero}
\par\medskip\hrule
\bigskip
\noindent{\large \textbf{Definition 5.} \textit{Intersecting with the empty set yields the empty set} \hfill Finset \cdot intersection \cdot empty is the right annihilator}\par
\smallskip\noindent\textit{Coordinate: Finset \cdot intersection \cdot empty is the right annihilator}\par
\smallskip\noindent\textit{Source: graham-knuth-patashnik}\par
\medskip
\noindent\textbf{Goal.} Intersecting with the empty set yields the empty set\par
\noindent$\displaystyle \forall s \in Finset\quad s ∩ empty = empty$\par
\medskip
\noindent
\begin{tabular}{@{} >{\raggedright\arraybackslash}p{0.06\textwidth} >{\raggedright\arraybackslash}p{0.47\textwidth} >{\raggedleft\arraybackslash}p{0.41\textwidth} @{}}
\textbf{\#} & \textbf{Proof} & \textbf{Mathematics} \\
\midrule
\textbf{1.} & We stipulate empty is the right annihilator for intersection on Finset — this is a definition, not a derived fact. & \\[0.45em]
\textbf{2.} & This follows directly from the definition: s ∩ empty equals empty. Hence proven. & $\displaystyle s ∩ empty = empty$ \\[0.45em]
\bottomrule
\end{tabular}
\label{thm:Finset-intersection-empty_is_the_right_annihilator}
\par\medskip\hrule
\bigskip
\noindent{\large \textbf{Definition 6.} \textit{Intersecting the empty set on the left yields the empty set} \hfill Finset \cdot intersection \cdot empty is the left annihilator}\par
\smallskip\noindent\textit{Coordinate: Finset \cdot intersection \cdot empty is the left annihilator}\par
\smallskip\noindent\textit{Source: graham-knuth-patashnik}\par
\medskip
\noindent\textbf{Goal.} Intersecting the empty set on the left yields the empty set\par
\noindent$\displaystyle \forall s \in Finset\quad empty ∩ s = empty$\par
\medskip
\noindent
\begin{tabular}{@{} >{\raggedright\arraybackslash}p{0.06\textwidth} >{\raggedright\arraybackslash}p{0.47\textwidth} >{\raggedleft\arraybackslash}p{0.41\textwidth} @{}}
\textbf{\#} & \textbf{Proof} & \textbf{Mathematics} \\
\midrule
\textbf{1.} & We stipulate empty is the left annihilator for intersection on Finset — this is a definition, not a derived fact. & \\[0.45em]
\textbf{2.} & This follows directly from the definition: empty ∩ s equals empty. Hence proven. & $\displaystyle empty ∩ s = empty$ \\[0.45em]
\bottomrule
\end{tabular}
\label{thm:Finset-intersection-empty_is_the_left_annihilator}
\par\medskip\hrule
\end{document}
